% Frontespizio del PDF (latex_elements["maketitle"] in conf.py)
\begin{titlepage}
\begin{tikzpicture}[remember picture,overlay]
  % fascia nera con filetto arancione
  \fill[metixnero] (current page.north west) rectangle ([yshift=-14cm]current page.north east);
  \fill[metixarancio] ([yshift=-14cm]current page.north west) rectangle ([yshift=-14.25cm]current page.north east);
  % motivo: colonna di layer che diventa celle poliedriche
  \begin{scope}[shift={([xshift=-5.4cm,yshift=-7.2cm]current page.north east)},scale=1.05,line width=0.7pt]
    \foreach \k in {0,1,2,3}{
      \draw[white,opacity=0.55] ({-2.2},{-1.6+0.32*\k*\k*0.6+0.25*\k}) -- (2.2,{-1.6+0.32*\k*\k*0.6+0.25*\k});
    }
    \foreach \x in {-2.2,-1.1,0,1.1,2.2}{\draw[white,opacity=0.55] (\x,-1.6) -- (\x,0.7);}
    \begin{scope}[metixarancio]
      \draw (-2.2,0.7) -- (-1.4,1.9) -- (-0.2,1.5) -- (0.6,2.4) -- (1.7,1.8) -- (2.2,0.7);
      \draw (-1.1,0.7) -- (-0.2,1.5); \draw (0,0.7) -- (-0.2,1.5); \draw (1.1,0.7) -- (1.7,1.8);
      \draw (-1.4,1.9) -- (-1.6,2.9) -- (-0.4,3.1) -- (0.6,2.4); \draw (-0.4,3.1) -- (-0.2,1.5);
      \draw (0.6,2.4) -- (1.4,3.2) -- (2.3,2.8) -- (1.7,1.8);
    \end{scope}
    \fill[metixarancio] (-2.2,-1.6) rectangle (2.2,-1.78);
  \end{scope}
\end{tikzpicture}
\vspace*{0.6cm}
{\metixsans\color{white!65!metixnero}\small\bfseries\addfontfeatures{LetterSpace=18}MANUALE TECNICO\par}
\vspace{10pt}
\includegraphics[height=1.75cm]{marchio-bianco.pdf}\par
\vspace{12pt}
{\metixsans\color{white}\LARGE Mesher poliedrico \emph{layer-first}\par}
\vspace{6pt}
{\metixsans\color{white!65!metixnero}\large Costruzione, algoritmi, origini\par}
\vspace{8.4cm}
\begin{minipage}{0.84\textwidth}
\large\color{metixinchiostro}
Questo documento descrive in dettaglio come MetiX costruisce una mesh poliedrica per OpenFOAM
a partire da una superficie triangolata: la preparazione della geometria, la parete, i layer di
strato limite, il core di Voronoi, l'assemblaggio duale, i controlli di qualità e gli sviluppi
sperimentali di novembre 2026. Per ogni metodo si indica se è ripreso dalla letteratura o se è
una costruzione nostra, e lo stato in cui si trova.
\end{minipage}
\vfill
{\metixsans\color{metixgrigio}\small Progetto MetiX \hfill Stato del codice: 13 novembre 2026 \hfill \texttt{metix-mesh/metix\_mesh/dualmesh}}
\end{titlepage}
